% R. Brown's translations of two letters to L. Breen, with his note and comments (cote 134-2): archive cote 134-2, site page 141 (PDF page 142).
% translation of the letter of 17-19 July 1975. Draft compared: ps1.tex, lines 1612-1621.
% AI-transcribed by Claude Opus 5.5 (claude-opus-5-5, adaptive thinking, effort high), Message Batches API, batch brown_full (msgbatch_01Md3i5yK1PxkpJn5GHfEExE), 2026-11-13; not yet checked by eye.
\scanpage{134-2/141}
\mspage{-14-}

limit) the knowledge of \uline{everything}.  And it depends in a convenient
sense, functorially on $X$ .  In the case of a scheme of finite
type on an algebraically closed  field $k$  say, the strongest
cohomological \uline{finiteness} \uline{theorem} in homotopy are expressed truly
precisely in terms of fine homotopy type, and say that the
ordinary homotopy types which are their elementary constituents
are essentially ``finite polyhedra'' - and even compact manifolds
with boundary - or in more precise fashion, the profinite
completions (in the sense of Artin-Mazur) prime to $p$ (the
characteristic) of such a space.  One sees clearly how to carry
out such a programme in characteristic  $0$ , but one has  more
fun, even    mystery , in the case  $p > 0$ , for the varieties
which, even  birationally, do not lift to characteristic $0$ !

From these essentially geometric thoughts, I have not at
this moment drawn up a precise programme for developing an
adequate algebraic structure to express them.  I restrict myself
to several marginal remarks.

For a long time I  have been intrigued by the idea of a
``linearisation'' of an (ordinary) homotopy type, i.e. questions of
the type: if  $X$  is a homotopy type, how much cohomological
information of the type: cohomology of  $X$  with variable co-%
efficients  $M$ (constant or twisted constant), multiplicative
structure  $H^{i}(X, M) \times H^{j}(X, N) \to H^{i+j}(X, M \otimes N)$ ,  then eventually
other cohomology operations - is it necessary to have to reconstruct
entirely the homotopy type (say, in this first setting, of the
type of derived pre-category, with the knowledge of the fundamental

\margin{142}
